\documentclass[12pt,a4paper,openright,twoside]{book}

% Cross-references into the book: \ref{ch06:prop:monotone} prints the
% book's own number. Build the book first (make book), then make solutions.
\usepackage{xr-hyper}
\usepackage{structecon}
\usepackage{subfiles}
\externaldocument{build/structural_economics_book}

% One heading per exercise: \exercisesolution{6.1}{Computing centralities}
\newcommand{\exercisesolution}[2]{%
  \section*{Exercise #1\quad #2}%
  \addcontentsline{toc}{section}{#1\quad #2}%
  \markright{Exercise #1}}
% Sub-part marker inside a solution: \part{a}
\newcommand{\solpart}[1]{\par\medskip\noindent\textbf{(#1)}\enspace}

\begin{document}

\frontmatter
\thispagestyle{empty}
\vspace*{3cm}
\begin{center}
{\fontsize{26}{32}\selectfont\bfseries STRUCTURAL ECONOMICS}\\[0.8em]
{\Large\itshape The General Case}\\[2.5cm]
{\LARGE Solutions Manual}\\[3cm]
{\Large Werner L.\ Hernani-Limarino}\\[4cm]
{\large DRAFT --- \today}
\end{center}
\clearpage

\chapter*{About this manual}
This manual gives detailed solutions to every exercise in \emph{Structural Economics: The General Case}.
Numbered references (Proposition~6.4, equation~(7.12)) point to the book.

\paragraph{Numbers.}
Every numerical answer is produced by a script in \texttt{code/solutions/solNN.py} (Python, \texttt{numpy}), run from the repository root: \texttt{python3 code/solutions/sol06.py}.
The scripts print each value used here, so an instructor can change a parameter and regenerate the answers.

\paragraph{Data exercises.}
Exercises that ask students to work with real data (INE, OECD, CEPAL, the Atlas of Economic Complexity) are answered in two parts.
First comes the method, step by step.
Then come the results, wherever the data could be fetched into the repository's \texttt{data/} directory, with the source and retrieval date.
Where the data could not be fetched, the manual says so and gives the method and the expected qualitative result only.

\paragraph{Open-ended questions.}
For interpretive questions, the manual gives a model answer: the argument a strong student should make, the formal object it rests on, and common mistakes.

\tableofcontents

\mainmatter
\IfFileExists{solutions/sol01.tex}{\subfile{solutions/sol01}}{}
\IfFileExists{solutions/sol02.tex}{\subfile{solutions/sol02}}{}
\IfFileExists{solutions/sol03.tex}{\subfile{solutions/sol03}}{}
\IfFileExists{solutions/sol04.tex}{\subfile{solutions/sol04}}{}
\IfFileExists{solutions/sol05.tex}{\subfile{solutions/sol05}}{}
\IfFileExists{solutions/sol06.tex}{\subfile{solutions/sol06}}{}
\IfFileExists{solutions/sol07.tex}{\subfile{solutions/sol07}}{}
\IfFileExists{solutions/sol08.tex}{\subfile{solutions/sol08}}{}
\IfFileExists{solutions/sol09.tex}{\subfile{solutions/sol09}}{}
\IfFileExists{solutions/sol10.tex}{\subfile{solutions/sol10}}{}
\IfFileExists{solutions/sol11.tex}{\subfile{solutions/sol11}}{}
\IfFileExists{solutions/sol12.tex}{\subfile{solutions/sol12}}{}
\IfFileExists{solutions/sol13.tex}{\subfile{solutions/sol13}}{}
\IfFileExists{solutions/sol14.tex}{\subfile{solutions/sol14}}{}
\IfFileExists{solutions/sol15.tex}{\subfile{solutions/sol15}}{}
\IfFileExists{solutions/sol16.tex}{\subfile{solutions/sol16}}{}
\IfFileExists{solutions/sol17.tex}{\subfile{solutions/sol17}}{}
\IfFileExists{solutions/sol18.tex}{\subfile{solutions/sol18}}{}
\IfFileExists{solutions/sol19.tex}{\subfile{solutions/sol19}}{}
\IfFileExists{solutions/sol20.tex}{\subfile{solutions/sol20}}{}
\IfFileExists{solutions/sol21.tex}{\subfile{solutions/sol21}}{}
\IfFileExists{solutions/sol22.tex}{\subfile{solutions/sol22}}{}
\IfFileExists{solutions/sol23.tex}{\subfile{solutions/sol23}}{}
\IfFileExists{solutions/sol24.tex}{\subfile{solutions/sol24}}{}
\IfFileExists{solutions/sol25.tex}{\subfile{solutions/sol25}}{}
\IfFileExists{solutions/sol26.tex}{\subfile{solutions/sol26}}{}
\IfFileExists{solutions/sol27.tex}{\subfile{solutions/sol27}}{}

\backmatter
\bibliography{references}

\end{document}
